\subsection*{并行与迭代}

\begin{table}[htbp]\centering\small
\begin{tabularx}{\linewidth}{@{}lX@{}}
\toprule
\textbf{语法} & \textbf{含义} \\
\midrule
\texttt{parallel p by N} & N 个独立实例，索引为 0 至 N-1 \\
\texttt{parallel px, py by [M, N]} & 同一并行层次中的两个坐标轴 \\
\texttt{parallel p by A, q by B} & 逗号分隔的多个并行层次 \newline{}（\texttt{p} 与 \texttt{q} 对应不同并行层次） \\
\texttt{parallel q by N : thread} & 映射到线程块内的线程 \\
\texttt{foreach index = m, n, k in [a, b, c]} & 在 \texttt{foreach} 中对具名元组解构 \\
\texttt{lhs.span(0)} & 取出张量形状的一个维度 \\
\texttt{p \# m} & 把外部分块索引 \texttt{p} 与内部偏移 \texttt{m} 合成为全局坐标 \\
\bottomrule
\end{tabularx}
\end{table}

2.12 节概述 MMA 的计算步骤，3.7 节给出完整示例。
